%! Simplifying Radicals; Adding \& Subtracting — Unit 5, Lesson 5.2 — Group Activity (Answer Key)
\documentclass[10pt]{article}
\usepackage{algebra2-article}
\usepackage{algebra2-key}   % pulls in -boxes; do NOT also load -boxes
\newcommand{\TallMath}[1]{$\displaystyle #1\rule[-1.4em]{0pt}{3.2em}$}
% In-formula answer slot (\ans is text-mode and must never appear inside $...$).
\newcommand{\mans}[1]{{\color{keyred}\mathbf{#1}}}

\begin{document}

\pageheader{Unit 5, Lesson 5.2}{Group Activity --- Answer Key}
\namepartnerperiod

\vspace{0.08in}

\begin{headlinebox}{forestbg}
\small\textbf{Hiding in Plain Sight.} Two radicals that look nothing alike can be the same radical
wearing different clothes --- and you will never know until you simplify them. With your partner,
simplify \emph{first} and decide \emph{second}. Say out loud which perfect power you pulled out and
why it was the largest one. Assume all variables represent \textbf{positive} real numbers.
\end{headlinebox}

\vspace{0.06in}

\begin{tcolorbox}[colback=white, colframe=black!40, title=\textbf{Tier R --- Remediate},
    fonttitle=\bfseries, arc=1.5mm, left=3mm, right=3mm, top=2mm, bottom=2mm, breakable]
\textbf{1. Find the largest perfect square, then simplify.} The first row is done for you.
\begin{center}
\renewcommand{\arraystretch}{1.7}
\begin{tabularx}{\linewidth}{@{}l X X X@{}}
\hline
\textbf{Radical} & \textbf{Largest perfect square factor} & \textbf{Rewrite} & \textbf{Simplified} \\
\hline
$\sqrt{18}$  & $9$ & $\sqrt{9\cdot 2}$ & $3\sqrt2$ \\
$\sqrt{32}$  & \ans{$16$} & \ans{$\sqrt{16\cdot 2}$} & \ans{$4\sqrt2$} \\
$\sqrt{44}$  & \ans{$4$}  & \ans{$\sqrt{4\cdot 11}$} & \ans{$2\sqrt{11}$} \\
$\sqrt{75}$  & \ans{$25$} & \ans{$\sqrt{25\cdot 3}$} & \ans{$5\sqrt3$} \\
$\sqrt{90}$  & \ans{$9$}  & \ans{$\sqrt{9\cdot 10}$} & \ans{$3\sqrt{10}$} \\
$\sqrt{128}$ & \ans{$64$} & \ans{$\sqrt{64\cdot 2}$} & \ans{$8\sqrt2$} \\
\hline
\end{tabularx}
\end{center}

\vspace{3pt}
\textbf{2. Finished, or not finished?} Mark each \textbf{D} (done --- already in simplest form) or
\textbf{N} (not done). If it is N, finish it.
\begin{center}
\renewcommand{\arraystretch}{1.7}
\begin{tabularx}{\linewidth}{@{}X X X@{}}
(a) \; $3\sqrt5$ \quad \ans{D} &
(b) \; $2\sqrt{12}$ \quad \ans{N} \; \ans{$4\sqrt3$} &
(c) \; $\sqrt{15}$ \quad \ans{D} \\
(d) \; $5\sqrt{8}$ \quad \ans{N} \; \ans{$10\sqrt2$} &
(e) \; $\sqrt[3]{9}$ \quad \ans{D} &
(f) \; $\sqrt[3]{24}$ \quad \ans{N} \; \ans{$2\sqrt[3]{3}$} \\
\end{tabularx}
\end{center}
{\small\emph{Watch (e):} $9=3^{2}$ is a perfect \emph{square}, but under a \emph{cube} root that
buys nothing --- there is no perfect cube inside $9$, so it is already done.}

\vspace{3pt}
\textbf{3. Find the families.} Simplify all six, then sort them into two groups of like radicals.
\begin{center}
\renewcommand{\arraystretch}{1.6}
\begin{tabularx}{\linewidth}{@{}X X X@{}}
$\sqrt{12}=$ \ans{$2\sqrt3$} & $\sqrt{50}=$ \ans{$5\sqrt2$} & $\sqrt{27}=$ \ans{$3\sqrt3$} \\
$\sqrt{8}=$ \ans{$2\sqrt2$} & $\sqrt{75}=$ \ans{$5\sqrt3$} & $\sqrt{32}=$ \ans{$4\sqrt2$} \\
\end{tabularx}
\end{center}
\begin{itemize}[leftmargin=1.6em, itemsep=3pt]
  \item The $\sqrt{3}$ family: \ans{$\sqrt{12},\ \sqrt{27},\ \sqrt{75}$} \; Their sum: \ans{$10\sqrt3$}
  \item The $\sqrt{2}$ family: \ans{$\sqrt{50},\ \sqrt{8},\ \sqrt{32}$} \; Their sum: \ans{$11\sqrt2$}
\end{itemize}

\vspace{2pt}
\textbf{4.} Why can the two sums above not be combined into a single term? \ansline{They are not
like radicals --- one has radicand $3$ and the other has radicand $2$. It is the same reason
$10a+11b$ stays as it is: there is nothing to collect. $10\sqrt3+11\sqrt2$ is already in simplest
form.}
\end{tcolorbox}

\begin{tcolorbox}[colback=white, colframe=black!40, title=\textbf{Tier A --- Approaching Proficiency},
    fonttitle=\bfseries, arc=1.5mm, left=3mm, right=3mm, top=2mm, bottom=2mm, breakable]
\textbf{1. Simplify completely.}
\begin{center}
\renewcommand{\arraystretch}{2.0}
\begin{tabularx}{\linewidth}{@{}X X X@{}}
(a) \; $\sqrt{162}=$ \ans{$9\sqrt2$} &
(b) \; $\sqrt[3]{128}=$ \ans{$4\sqrt[3]{2}$} &
(c) \; $\sqrt{48x^{7}}=$ \ans{$4x^{3}\sqrt{3x}$} \\
(d) \; $\sqrt[3]{81y^{8}}=$ \ans{$3y^{2}\sqrt[3]{3y^{2}}$} &
(e) \; $\sqrt{\dfrac{45}{4}}=$ \ans{$\dfrac{3\sqrt5}{2}$} &
(f) \; $\dfrac{\sqrt{98}}{\sqrt{2}}=$ \ans{$7$} \\
\end{tabularx}
\end{center}
{\small\emph{Work:} (a) $162=81\cdot2$. \; (b) $128=64\cdot2$. \; (c) $48=16\cdot3$ and
$x^{7}=x^{6}\cdot x$. \; (d) $81=27\cdot3$ and $y^{8}=y^{6}\cdot y^{2}$ (divide $8$ by $3$:
quotient $2$, remainder $2$). \; (e) $\sqrt{45}/\sqrt4=3\sqrt5/2$. \; (f) run the quotient property
backwards: $\sqrt{98/2}=\sqrt{49}$.}

\vspace{3pt}
\textbf{2. Combine.} Simplify every term \emph{first}, then collect like radicals.
\begin{center}
\renewcommand{\arraystretch}{2.0}
\begin{tabularx}{\linewidth}{@{}X X@{}}
(a) \; $2\sqrt{18}+3\sqrt{50}=$ \ans{$21\sqrt2$} &
(b) \; $\sqrt{80}-\sqrt{45}+\sqrt{20}=$ \ans{$3\sqrt5$} \\
(c) \; $\sqrt{27x}+\sqrt{12x}=$ \ans{$5\sqrt{3x}$} &
(d) \; $4\sqrt{24}-\sqrt{54}+\sqrt{6}=$ \ans{$6\sqrt6$} \\
\end{tabularx}
\end{center}
{\small\emph{Work:} (a) $6\sqrt2+15\sqrt2$. \; (b) $4\sqrt5-3\sqrt5+2\sqrt5$. \;
(c) $3\sqrt{3x}+2\sqrt{3x}$. \; (d) $8\sqrt6-3\sqrt6+\sqrt6$.}

\vspace{3pt}
\textbf{3. Find the mistake.} Each student made exactly one error. Name it and give the correct
answer.
\begin{center}
\renewcommand{\arraystretch}{1.7}
\begin{tabularx}{\linewidth}{@{}l X X@{}}
\hline
\textbf{Student work} & \textbf{What went wrong} & \textbf{Correct answer} \\
\hline
$\sqrt{20}+\sqrt{5}=\sqrt{25}=5$ & \ans{split the radical over addition --- there is no such rule} & \ans{$3\sqrt5$} \\
$3\sqrt7+2\sqrt7=5\sqrt{14}$ & \ans{added the radicands too; only the numbers in front add} & \ans{$5\sqrt7$} \\
$\sqrt[3]{54}=\sqrt[3]{9\cdot 6}=3\sqrt[3]{6}$ & \ans{used a perfect \emph{square} under a \emph{cube} root; $9$ does not come out} & \ans{$3\sqrt[3]{2}$} \\
\hline
\end{tabularx}
\end{center}
{\small The first student's own numbers refute them: $\sqrt{20}+\sqrt5\approx 4.47+2.24=6.71$,
not $5$. The third student happened to write ``$3$'' outside, which looks right --- but they got it
from $\sqrt{9}$, not from $\sqrt[3]{27}$; $54=27\cdot 2$.}

\vspace{3pt}
\textbf{4. Exact beats approximate.} A rectangular garden plot measures $\sqrt{75}$ meters by
$\sqrt{12}$ meters.
\begin{enumerate}[label=(\alph*), leftmargin=1.9em, itemsep=4pt]
  \item Write each side in simplest radical form: \; length \ans{$5\sqrt3$}, \; width \ans{$2\sqrt3$}
  \item Find the \textbf{exact} perimeter as a single simplified radical: \ans{$14\sqrt3$ m}
  \item Now approximate it to the nearest tenth of a meter: \ans{$\approx 24.2$ m}
  \item Why would a builder want the answer in part (b) and not just the one in part (c)?
        \ansline{$14\sqrt3$ is exact --- no rounding has happened yet, so it can be scaled, doubled,
        or fed into another calculation without the error growing. $24.2$ has already thrown
        information away, and every later step inherits that loss. Round \emph{once}, at the end.}
\end{enumerate}
\end{tcolorbox}

\begin{tcolorbox}[colback=white, colframe=black!40, title=\textbf{Tier E --- Extension},
    fonttitle=\bfseries, arc=1.5mm, left=3mm, right=3mm, top=2mm, bottom=2mm, breakable]
\textbf{Part 1 --- Prove the error is never a lucky guess.} Everyone knows
$\sqrt{a}+\sqrt{b}\ne\sqrt{a+b}$ ``usually.'' You are going to prove it is \emph{never} true unless
one of the two numbers is zero. Let $a\ge 0$ and $b\ge 0$, and \emph{suppose} for a moment that
\[
  \sqrt{a}+\sqrt{b}=\sqrt{a+b}.
\]
\begin{enumerate}[label=(\alph*), leftmargin=1.9em, itemsep=5pt]
  \item Square both sides. The left side is a binomial squared --- expand it carefully.
        \ansline{$\left(\sqrt a+\sqrt b\right)^{2}=a+2\sqrt{ab}+b$, and
        $\left(\sqrt{a+b}\right)^{2}=a+b$, so $a+2\sqrt{ab}+b=a+b$.}
  \item Subtract $a+b$ from both sides. What is left? \ans{$2\sqrt{ab}=0$}
  \item Solve that equation. What must be true about $a$ and $b$? \ans{$ab=0$, so $a=0$ or $b=0$}
  \item State the conclusion in a sentence, then check it against $a=9$, $b=16$. \ansline{The two
        sides are equal \textbf{only} when at least one of the numbers is zero; for any two
        \emph{positive} numbers the equation is false, always. Check: $\sqrt9+\sqrt{16}=7$ while
        $\sqrt{25}=5$, and neither $9$ nor $16$ is zero --- so of course they disagree.\\[3pt]
        The extra term $2\sqrt{ab}$ is exactly what the mistake throws away.}
\end{enumerate}

\vspace{5pt}
\textbf{Part 2 --- A staircase of radicals.} Simplify each term, then add.
\[
  \sqrt{2}+\sqrt{8}+\sqrt{18}+\sqrt{32}+\sqrt{50} \;=\;
  \mans{\sqrt2}+\mans{2\sqrt2}+\mans{3\sqrt2}+\mans{4\sqrt2}+\mans{5\sqrt2}
  \;=\; \mans{15\sqrt2}
\]
\begin{enumerate}[label=(\alph*), leftmargin=1.9em, itemsep=5pt]
  \item Look at the radicands $2,\,8,\,18,\,32,\,50$. Write each one in the form
        $2\cdot(\text{something})^{2}$: \ans{$2\cdot1^{2},\ 2\cdot2^{2},\ 2\cdot3^{2},\ 2\cdot4^{2},\ 2\cdot5^{2}$}
  \item Explain in one sentence why $\sqrt{2k^{2}}=k\sqrt{2}$. \ansline{By the product property,
        $\sqrt{2k^{2}}=\sqrt{k^{2}}\cdot\sqrt2=k\sqrt2$ for $k>0$ --- the perfect square $k^{2}$
        comes out and the $2$ has no partner, so it stays.}
  \item What is the \emph{next} term in the staircase, and what does it simplify to?
        \ans{$\sqrt{72}=6\sqrt2$ \;(since $72=2\cdot 6^{2}$)}
  \item Predict the total of the first $n$ terms, in terms of $n$: \ans{$\dfrac{n(n+1)}{2}\sqrt2$} \\
        \emph{Hint:} you added $1+2+3+\dots$ Check your formula against the $n=5$ answer above.
        \ansline{Every term contributes its position number, so the total is
        $(1+2+\dots+n)\sqrt2=\frac{n(n+1)}{2}\sqrt2$. At $n=5$: $\frac{5\cdot6}{2}\sqrt2=15\sqrt2$,
        which matches.}
  \item Build the same staircase with $3$ in place of $2$: simplify
        $\sqrt{3}+\sqrt{12}+\sqrt{27}+\sqrt{48}=$ \ans{$10\sqrt3$}
\end{enumerate}

\vspace{5pt}
\textbf{Part 3 --- Two functions that were secretly the same shape.} Consider
$f(x)=\sqrt{8x}$ and $g(x)=\sqrt{2x}$.
\begin{center}
\renewcommand{\arraystretch}{1.3}
\begin{tabular}{|c|c|c|c|}
\hline
$x$ & $2$ & $8$ & $18$ \\ \hline
$g(x)=\sqrt{2x}$ & \ans{$2$} & \ans{$4$} & \ans{$6$} \\ \hline
$f(x)=\sqrt{8x}$ & \ans{$4$} & \ans{$8$} & \ans{$12$} \\ \hline
\end{tabular}
\end{center}
\begin{enumerate}[label=(\alph*), leftmargin=1.9em, itemsep=5pt]
  \item In every column, $f(x)$ is exactly \ans{$2$} times $g(x)$.
  \item Prove it with the product property: $\sqrt{8x}=\sqrt{\mans{4}\cdot 2x}=$ \ans{$2\sqrt{2x}$}
  \item So $f$ is not a new function at all --- it is $g$ stretched \textbf{vertically} by a factor
        of \ans{$2$}. Which of the two do $f$ and $g$ share: the same domain, the same
        endpoint, or both? \ansline{Both. A vertical stretch multiplies every output by $2$ but
        touches no input, so the domain is still $x\ge0$; and the endpoint $(0,0)$ has output $0$,
        and $2\cdot 0=0$, so it does not move either. Only the steepness changed.}
\end{enumerate}
\end{tcolorbox}

\begin{teachernote}
\textbf{Tier R} (under twelve minutes). Item 2(e) is the one to watch: most students will call
$\sqrt[3]{9}$ ``not done'' because they see a $9$ and think ``perfect square.'' Worth a whole-class
share --- the \emph{index} decides what counts. In item 3, do not let groups sort before they
simplify; the sorting is impossible otherwise, which is the point.

\textbf{Tier A} item 3 is the highest-value item on the page. Have three groups present the three
errors: the first student's own decimals ($4.47+2.24=6.71\ne5$) refute them, and the third
student's answer \emph{looks} right, which is why it teaches. Item 4(d) names a modeling habit ---
keep the exact form, round once, at the end.

\textbf{Tier E} Part 1 is a real proof, and accessible: it all turns on expanding a binomial
square. The punch line is that $2\sqrt{ab}$ is exactly the term the classic error deletes. Part 2
closes with a genuine generalization; groups that get (d) can be asked whether \emph{every}
staircase $\sqrt{c}+\sqrt{4c}+\sqrt{9c}+\dots$ behaves the same way (it does). Part 3 is a bridge to
Lesson 5.4 --- simplifying $\sqrt{8x}$ into $2\sqrt{2x}$ reveals a \textbf{vertical stretch}, the
same $a$ in $y=a\sqrt{x-h}+k$ students meet in two days.
\end{teachernote}

\end{document}
